\section{Undecidability with functional inputs}\label{sec:functional-hardness}

The same controls that give the ground-input decision theorem support
an undecidable comparison problem once free variables of type $o\to o$
are allowed. The reduction uses the growth estimates, not any modular
property of the control. The final corollary records the same negative
result for the unchanged standard control.

The published undecidable problem used here is solvability of a finite
system of linear equations
\[
 \sum_i a_{ji}(X)p_i(X)-g_j(X)=0,
 \qquad a_{ji},g_j\in\mathbb Z[X],\quad p_i\in\mathbb N[X].
\]
It is linear in the unknown polynomials, not a system of arbitrary
Diophantine products of unknowns.  Narendran's theorem gives this precise polynomial-semiring boundary
\cite{narendran}.
Zero polynomials are allowed.  We use undecidability, without importing a
stronger many-one completeness classification of that source problem.

For ground sources define the eraser $e(U,V)=(\lambda u:o.U)V$,
with a fresh unused binder. Its literal scalar is
\[
 e(U,V)_0=g(3U_0+V_0+2),\qquad g(n)=\psi_f(\omega+n)=f(n+2)+1.
\]
The function $g$ is strictly increasing with gaps at least two.

Put $\tau=o\to o$.  For actual sources $A,B:\tau$ define
\begin{align*}
 \mathsf P(A,B)&=(\lambda z:\tau.A)B,\qquad
 \mathsf H(A)=R^o0(\lambda n:o.A),\\
 \mathsf S(A,B)&=(\lambda z:(o\to\tau).A)(\lambda n:o.B).
\end{align*}
The original argument order is
$R^\rho:o\to(o\to\rho\to\rho)\to\rho\to\rho$:
its first argument $0:o$ is the recursion test, not its base value.
Thus $R^\rho0(\lambda n:o.A)$ has type $\rho\to\rho$
when $A:\rho\to\rho$, also at the higher type used below.
The displayed bound names are unused and fresh. These literal identities
are uniform over admissible parsing policies: their unused-binder partials
do not inspect a rebracketing inside an active variable occurrence.  Writing
$a_i=\can A_i$, $b_i=\can B_i$, the literal canonical clauses give
\begin{align*}
 \can{\mathsf P(A,B)}_1&=8a_1+4b_1+8,\\
 \can{\mathsf S(A,B)}_1&=8a_1+8b_1+16,\\
 \can{\mathsf S(A,B)}_0
   &=\psi_f\bigl(\omega\otimes(8a_1+8b_1+16)+9a_0+5b_0+13\bigr),\\
 \can{\mathsf H(A)}_1&=2\omega\otimes a_1+3\omega.
\end{align*}
For example the unused $\tau$ binder has vector
$(9a_0+8,2a_1+2,2)$, and its box with $B$ gives the first formula.
For the second binder its scalar is again $9a_0+8$ and its first field
$2a_1+2$; the ground lambda on $B$ has vector
$(5b_0+4,2b_1+2)$, giving both $\mathsf S$ formulas.
The recursor prime followed by its two displayed arguments gives
the $\mathsf H$ formula.  These calculations use actual source syntax;
they do not introduce an arbitrary ordinal-valued primitive.

\begin{lemma}[Signed polynomial pairs]\label{lem:pair-compiler}
Every signed linear expression $D$ in ordinal unknowns with coefficients
in $\mathbb Z[\omega]$ has effectively constructed actual $A_D,B_D:\tau$
whose first-coordinate difference is $2^kD$ for a finite $k$.
All free functional variables have type $\tau$, and internal types have
level at most two.  Differences are taken only in the signed Cantor-monomial
group; each of the two source coordinates remains a positive ordinal.
\end{lemma}
\begin{proof}
Let $C=\lambda u:o.u$, with $C_1=1$.  The first fields of
$\mathsf P(C,C)$ and $\mathsf S(C,C)$ are 20 and 32.  Take sixteen pairs:
eight copies of $(\mathsf S(C,C),\mathsf P(C,C))$ of difference 12,
five of $(C,\mathsf P(C,C))$ of difference $-19$, and three of $(C,C)$.
Combine them componentwise in a balanced $\mathsf S$ tree of height four.
Constants cancel in every paired combination, so the resulting difference is
\[
 8^4(8\cdot12-5\cdot19)=2^{12}.
\]
For a free $X:\tau$, the pair
$(\mathsf P(X,C),\mathsf P(C,C))$ has difference $8(x_1-1)$.
Three multiplications by eight followed by a paired $\mathsf S$ combination
with the constant pair give $2^{15}x_1$.

The pair operation $(A,B)\mapsto(\mathsf S(A,C),\mathsf S(B,C))$
multiplies a difference by eight, and
$(A,B)\mapsto(\mathsf S(A,A),\mathsf S(B,B))$ multiplies it by sixteen.
Applying $\mathsf H$ to both entries multiplies a difference by $2\omega$;
swapping entries negates it.  Repeated addition constructs every finite
integer coefficient.  To add pairs with scale exponents $k_1,k_2$, choose
$K\ge\max(k_1,k_2)+6$.  Every integer at least six is $3u+4v$ with
$u,v\ge0$, so the two scale operations align both differences at $2^K$.
A paired $\mathsf S$ then gives their sum with scale $2^{K+3}$.
Induction on the finitely many signed monomials proves the assertion.
The displayed source combinators introduce no internal level above two.
\end{proof}

The two scalar coordinates of an equation pair need not agree when their
positive coordinates agree.  We therefore construct that alignment explicitly.

\begin{lemma}[Alignment and a full-vector detector]\label{lem:functional-detector}
At each fixed admissible valuation of the original inputs, the condition
$\can A_1=\can B_1$ is equivalent to the existence of two fresh
positive ground inputs making the full canonical vectors of
\[
 A'=(\lambda t:o.A)z,\qquad B'=(\lambda t:o.B)w
\]
equal.  Full-vector equality of any actual pair of type $\tau$ can in turn
be detected by ground sources $L,R$ satisfying $L_0\ge R_0$, with equality
exactly at vector equality and otherwise gap at least two.
\end{lemma}
\begin{proof}
The aligned first fields are $2a_1+2$ and $2b_1+2$, and their scalars are
\[
 \psi_f(\omega\otimes(2a_1+2)+5a_0+4+z_0),\qquad
 \psi_f(\omega\otimes(2b_1+2)+5b_0+4+w_0).
\]
Equal vectors imply $a_1=b_1$.  If this equality holds, choose
$z_0=5b_0+1$, $w_0=5a_0+1$.  Both tails and both vectors then agree.
These choices are made after the original valuation, so are not circular.

For the detector, now write $A,B$ for the pair being tested, and form
\[
 X=\mathsf S(A,A),\quad Y=\mathsf S(B,B),\quad
 Z=\mathsf S(A,B),\quad W=\mathsf S(B,A).
\]
Make the actual shared ground payload
$Q=W(Z(Y(X0)))$, and put $q=Q_0$.
The unconditional finite-tail lower bound for $\psi_f$ gives
\[
 q\ge1+x_0+y_0+z_0+w_0\ge5.
\]
Set $g_X=XQ$, and similarly $g_Y,g_Z,g_W$, and finally set
$L=e(g_X,g_Y)$, $R=e(g_Z,g_W)$.

Suppose first $a_1>b_1$.  Then $x_1>z_1=w_1$.  The ordinals
\[
 \alpha=\omega\otimes z_1+z_0+q+2,
 \qquad \beta=\omega\otimes x_1+x_0+q
\]
satisfy $\alpha<\beta$.  We must check a norm premise before comparing
their collapses.  Bounded norms of actual assignments and the published
natural-product norm inequality \cite[Propositions 2.2--2.3]{ww} give
\[
 \norm(\alpha)\le2\norm(z_1)+z_0+q+2
 \le3z_0+q+2\le4q+2
 \le2^{q+1}\le f(\norm(\beta)).
\]
Controlled strictness therefore gives $\psi_f(\beta)>\psi_f(\alpha)$.
Two consecutive fixed-limit tail-growth steps from Lemma~\ref{lem:tail-growth}
yield $g_{X,0}>4g_{Z,0}$.  The same calculation with $W$ gives
$g_{X,0}>4g_{W,0}$.  If $b_1>a_1$, use $Y$ in place of $X$.

If $a_1=b_1$ but $a_0\ne b_0$, the four $\mathsf S$ heads coincide.
Their scalar tails are
\[
 14a_0+13,\quad14b_0+13,\quad
 9a_0+5b_0+13,\quad9b_0+5a_0+13.
\]
The larger diagonal tail exceeds both cross tails by at least five.
Strict increase gives at least a two-unit scalar gap, so two fixed-tail
growth steps in the subsequent applications give again a diagonal $g$
value greater than four times each cross value.

In either unequal case,
$3g_{X,0}+g_{Y,0}>3g_{Z,0}+g_{W,0}$: even when the larger diagonal
is the one with coefficient one, its fourfold domination suffices.
The final eraser is increasing and has gaps at least two.
Equal input vectors give four equal $\mathsf S$ vectors, hence equality.
This proves the claimed detector on every admissible valuation, not only
on the specially chosen alignment witnesses.
\end{proof}

\begin{theorem}[Functional-input undecidability]
\label{thm:functional-hardness}
Full universal canonical weak comparison and strict comparison are
undecidable already for ground output, free-variable levels at most one,
and internal levels at most two.  The reduction at this level uses finitely
many comparison queries.  A computable single-query reduction from the
same polynomial-system unsatisfiability problem is available with internal
level at most three.
\end{theorem}
\begin{proof}
Let $\pi$ retain exactly the Cantor monomials with finite exponents.
It maps ordinals below $\varepsilon_0$ to $\mathbb N[\omega]$.
It preserves natural sum and multiplication: a natural sum of exponents
is finite exactly when both exponents are finite.  It consequently extends
to the signed Cantor-monomial group and respects multiplication by fixed
$\mathbb Z[\omega]$ coefficients.

For the level-two assertion, enumerate the finitely many possible sets of
zero polynomial unknowns, substitute zero for those, and require the others
to be positive ordinal first coordinates.  Compile the resulting equations
by Lemma~\ref{lem:pair-compiler}, align them separately, and apply
Lemma~\ref{lem:functional-detector}.  Aggregate the nonnegative ground
detectors in identically shaped eraser trees.  Root equality is exactly
simultaneous equation satisfaction.  Its absence is one universal strict
comparison; the conversion
\[
 \can{KR}_0\le\can{IL}_0\quad\Longleftrightarrow\quad
 g(5+R_0)\le g(3+L_0)\quad\Longleftrightarrow\quad L_0-R_0\ge2
\]
 gives the same weak query. Here $K=\lambda z:o.0$ and
$I=\lambda z:o.z$, with canonical vectors $(5,1)$ and $(3,1)$.
An empty polynomial system is first augmented by the tautological equation
$0=0$, so the same unsatisfiability-query convention applies.

A polynomial solution lifts in the matching zero pattern: each remaining
nonzero polynomial is an admissible positive first field, after choosing its
zeroth field sufficiently large for its norm.  Conversely any positive-ordinal
solution projects by $\pi$ to a polynomial solution of the original system.
Some projected unknowns may become zero, which is permitted in that system.
Thus an algorithm for either comparison would decide polynomial solvability
using the finitely many patterns, contradicting the published theorem.

For a single query choose instead a fixed positive ordinal $\Omega$ with
$\pi(\Omega)=0$, and transform each original equation to
\[
 \sum_i a_i(\omega)x_i-g(\omega)
                    -\left(\sum_i a_i(\omega)\right)\Omega=0.
\]
A polynomial solution $u_i$ lifts by $x_i=u_i\oplus\Omega>0$.
Conversely projection of any positive-ordinal solution solves the original
equations.  The transformation preserves simultaneous use of each unknown.

We must realize the extra constant in actual source syntax.  Put
$\sigma=\tau\to\tau$, $C_2=\lambda u:\tau.u$, and
$\mathsf H_2(A)=R^\tau0(\lambda n:o.A)$.  The canonical clauses give
$(C_2)_2=2$ and $(\mathsf H_2(A))_2=\omega+2\omega\otimes A_2$.
The closed $H=\mathsf H_2(\mathsf H_2(C_2))$ has
\[
 H_2=10\omega^2+\omega,
 \qquad\Omega=2^{10\omega^2+\omega}=\omega^{10\omega+1},
 \qquad\pi(\Omega)=0.
\]
Apply this same actual $H$ to the two entries of the constant pair of
Lemma~\ref{lem:pair-compiler}.  The box gives
$\can{HA}_1=\Omega\otimes(H_1+A_1)$, so their difference is
$2^{12}\Omega$.  Add this as a second closed seed to the pair compiler.
Only fixed $\omega$ shifts and signed sums of that seed are needed.
Compile the one transformed system and apply the same alignment and detector.
The closed $R^\tau$ raises the internal level bound to three, while every
free variable still has level at most one.  This is the promised single-query
reduction.
\end{proof}

\begin{corollary}[Functional inputs at the unchanged standard control]
\label{cor:standard-functional-hardness}
The weak and strict undecidability statements of
Theorem~\ref{thm:functional-hardness} also hold for the unchanged control
\(G(n)=F_5(n+100)\), with the same finite-query internal level bound two
and single-query internal level bound three.
\end{corollary}
\begin{proof}
The canonical-source interface requires a strictly increasing control
majorizing \(G\), and the numerical separation used in the detector
requires the adjacent growth estimate \eqref{eq:growth}. The standard
control meets these premises. Indeed, with \(A=G(n)\), the hierarchy
calculation in Lemma~\ref{lem:control} gives
\[
 G(n+1)\ge F_4(A+1)\ge F_1(A+1)
             \ge 2^{2^{A+1}}\ge2^{2^{G(n)}}.
\]
It is primitive recursive and has the same eventual finite-level bound.
No count-residue or modular-potential statement occurs in the source
compiler, alignment, detector or undecidability reduction. Their literal
constructions and proofs therefore apply with \(f=G\) throughout.
\end{proof}

